\documentclass[12pt,a4paper]{article}

\usepackage[utf8]{inputenc}
\usepackage[T1]{fontenc}
\usepackage[brazil]{babel}
\usepackage{geometry}
\usepackage{amsmath}
\usepackage{booktabs}
\usepackage{array}
\usepackage{float}
\usepackage{setspace}

\geometry{
    left=3cm,
    right=2cm,
    top=3cm,
    bottom=2cm
}

\onehalfspacing

\title{Exemplo - Análise de Soluções e Frequências no Problema das P-Medianas}
\author{Carlos Vinicius Fialho Rabelo}
\date{2026}

\begin{document}

\maketitle

\section{Item 1 -- Cálculo dos Custos das Soluções}

Para avaliar as soluções, considera-se a menor distância
entre cada cliente e as instalações selecionadas.
O custo total de uma solução corresponde à soma dessas
menores distâncias.

As soluções analisadas são
$\{0,7,12\}$, $\{1,7,12\}$ e $\{1,4,9\}$.

\subsection{Solução $\{0,7,12\}$}

\begin{itemize}
    \item Cliente 0: $|0-0|=0$, $|0-7|=7$,
    $|0-12|=12$. Menor distância: 0.
    \item Cliente 1: $1$, $6$, $11$.
    Menor distância: 1.
    \item Cliente 4: $4$, $3$, $8$.
    Menor distância: 3.
    \item Cliente 7: $7$, $0$, $5$.
    Menor distância: 0.
    \item Cliente 9: $9$, $2$, $3$.
    Menor distância: 2.
    \item Cliente 12: $12$, $5$, $0$.
    Menor distância: 0.
\end{itemize}

O custo total é:

\[
C(\{0,7,12\}) = 0+1+3+0+2+0 = 6
\]

\subsection{Solução $\{1,7,12\}$}

\begin{itemize}
    \item Cliente 0: $1$, $7$, $12$.
    Menor distância: 1.
    \item Cliente 1: $0$, $6$, $11$.
    Menor distância: 0.
    \item Cliente 4: $3$, $3$, $8$.
    Menor distância: 3.
    \item Cliente 7: $6$, $0$, $5$.
    Menor distância: 0.
    \item Cliente 9: $8$, $2$, $3$.
    Menor distância: 2.
    \item Cliente 12: $11$, $5$, $0$.
    Menor distância: 0.
\end{itemize}

O custo total é:

\[
C(\{1,7,12\}) = 1+0+3+0+2+0 = 6
\]

\subsection{Solução $\{1,4,9\}$}

\begin{itemize}
    \item Cliente 0: $1$, $4$, $9$.
    Menor distância: 1.
    \item Cliente 1: $0$, $3$, $8$.
    Menor distância: 0.
    \item Cliente 4: $3$, $0$, $5$.
    Menor distância: 0.
    \item Cliente 7: $6$, $3$, $2$.
    Menor distância: 2.
    \item Cliente 9: $8$, $5$, $0$.
    Menor distância: 0.
    \item Cliente 12: $11$, $8$, $3$.
    Menor distância: 3.
\end{itemize}

O custo total é:

\[
C(\{1,4,9\}) = 1+0+0+2+0+3 = 6
\]

\subsection{Resumo dos resultados}

\begin{table}[H]
\centering
\caption{Custos das soluções analisadas}
\begin{tabular}{cc}
\toprule
\textbf{Solução} & \textbf{Custo total} \\
\midrule
$\{0,7,12\}$ & 6 \\
$\{1,7,12\}$ & 6 \\
$\{1,4,9\}$ & 6 \\
\bottomrule
\end{tabular}
\end{table}

A análise mostra que as três soluções possuem o mesmo
custo total, igual a 6. Portanto, considerando apenas
a função de custo, não é possível afirmar que uma delas
seja melhor que as demais. As soluções estão empatadas,
não havendo vantagem em escolher uma como solução-base.

\section{Item 2 -- Cálculo das Frequências}

O cálculo das frequências consiste em contar quantas
vezes cada instalação aparece nas soluções analisadas.
Como cada solução contém três instalações e existem
três soluções, a quantidade total de ocorrências é:

\[
P \times N = 3 \times 3 = 9
\]

em que $P$ representa a quantidade de instalações e $N$ representa a quantidade de soluções.

\subsection{Frequência das instalações}

\begin{table}[H]
\centering
\caption{Frequência de ocorrência das instalações}
\begin{tabular}{cc}
\toprule
\textbf{Instalação} & \textbf{Quantidade nas soluções} \\
\midrule
$\{0\}$  & 1 \\
$\{1\}$  & 2 \\
$\{4\}$  & 1 \\
$\{7\}$  & 2 \\
$\{9\}$  & 1 \\
$\{12\}$ & 2 \\
\midrule
\textbf{Total} & \textbf{9} \\
\bottomrule
\end{tabular}
\end{table}

A soma das frequências é igual a 9, confirmando que
todas as ocorrências de instalações nas três soluções
foram contabilizadas.

\subsection{Frequência dos pares}

Além das instalações individuais, podem ser analisados
os pares de instalações presentes em cada solução.
Como cada solução contém três instalações, ela gera
três pares distintos.

\begin{table}[H]
\centering
\caption{Pares presentes em cada solução}
\begin{tabular}{cc}
\toprule
\textbf{Solução} & \textbf{Pares} \\
\midrule
$\{0,7,12\}$ &
$\{0,7\}$, $\{0,12\}$, $\{7,12\}$ \\
$\{1,7,12\}$ &
$\{1,7\}$, $\{1,12\}$, $\{7,12\}$ \\
$\{1,4,9\}$ &
$\{1,4\}$, $\{1,9\}$, $\{4,9\}$ \\
\bottomrule
\end{tabular}
\end{table}

O par $\{7,12\}$ aparece duas vezes, enquanto os demais
pares aparecem uma única vez. Dessa forma, o par
$\{7,12\}$ é o mais frequente entre as soluções
analisadas e pode ser utilizado como ponto de partida
para a construção de uma nova solução.

\section{Item 3 -- Utilização das Frequências na Construção
de Novas Soluções}

As frequências obtidas indicam que as instalações 1,
7 e 12 aparecem duas vezes nas soluções analisadas.
Além disso, o par $\{7,12\}$ também aparece duas vezes,
sendo o par mais frequente.

Essas informações podem ser utilizadas para orientar
a construção de novas soluções. Uma possibilidade é
usar as frequências como pesos em um sorteio,
atribuindo maior probabilidade de escolha às instalações
e aos pares que aparecem com maior frequência.

Essa estratégia mantém o componente aleatório da
construção, mas utiliza as informações extraídas das
soluções anteriores para orientar as escolhas.

Nesse contexto, o pool de soluções de elite funciona
como uma memória do algoritmo, armazenando soluções
consideradas boas. As frequências representam informações
extraídas dessa memória.

Por sua vez, o \textit{path-relinking} pode combinar
características de diferentes soluções de elite,
explorando caminhos entre elas em busca de soluções
de melhor qualidade.

Uma analogia possível é a consulta a um banco de dados:
as soluções armazenadas são observadas para identificar
padrões que podem ser utilizados na
construção de novas soluções. Entretanto, diferente de 
uma simples verificação de integridade de dados,
o objetivo do algoritmo é aproveitar essas informações
para explorar o espaço de soluções e buscar custos
menores.

\end{document}
